\chapter{Austrian Attack with 5... c5: Early Deviations}

Below are some typical positions found in the Pirc.

\section{Positional Plans}

\newchessgame[id=AustrianAttack1,
initialfigure=\thefigurecnt]
\mainline{1.e4 d6 2. d4 Nf6 3.Nc3 g6 4. f4} 

\chessboard[showmover=true]
\lfigure

This is the defining position for the Austrian Attack. Black plans to 
challenge White's central pawns with properly prepared 
c5 or e5. White plans is to move e5 or d5 (advance the central pawns).

The typical resulting positions have the following pawn structures:

\newchessgame[id=PawnStructureIsolani,
initialfigure=\thefigurecnt,
moveid=16b,
setwhite={pa2,pb2,pc2,pe4,pg2,ph2},
addblack={pa7,pb7,pc7,pf7,pg6,ph7},
storefen=PSIsolani]
\chessboard[inverse,showmover=false]
\lfigure

\newchessgame[id=PawnStructureBenoniPirc,
initialfigure=\thefigurecnt,
moveid=16b,
setwhite={pa2,pb2,pc2,pd5,pe4,pf4,pg2,ph2},
addblack={pa7,pb7,pc5,pd6,pe7,pf7,pg6,ph7},
storefen=PSBenoniPirc]
\chessboard[inverse,showmover=false]
\lfigure

\newchessgame[id=PawnStructureCentralWedge,
initialfigure=\thefigurecnt,
moveid=16b,
setwhite={pa2,pb2,pc2,pe5,pf4,pg2,ph2},
addblack={pa7,pb7,pc5,pe7,pf7,pg6,ph7},
storefen=PSCentralWedge]
\chessboard[inverse,showmover=false]
\lfigure

\newchessgame[id=PawnStructureDRam,
initialfigure=\thefigurecnt,
moveid=16b,
setwhite={pa2,pb2,pc2,pd5,pe4,pf4,pg2,ph2},
addblack={pa7,pb7,pc6,pd6,pe7,pf7,pg6,ph7},
storefen=PSDRam]
\chessboard[inverse,showmover=false]
\lfigure

\section{Early deviations   5.Bd3 Nc6 }

\newchessgame[id=EarlyDeviationI]
\mainline{1.e4 d6 2. d4 Nf6 3.Nc3 g6 4. f4 Bg7
5. Bd3 Nc6} 
\chessboard[inverse,showmover=false]
\lfigure

\xskakcomment{7.Bd3 Ng4 loses a pawm. 7.d5 Nb4 with Black's idea of
...c6 with the possibility of ...Nh6 ...Nc7. Note that ...c6 permits Black's
queen to gain access to the a5-d8 diagonal. ...Qb6 permits ideal access to 
White's weak f2 square.}

\section{Key Positions}

\newchessgame[id=EarlyDeviation]
\mainline{1.e4 d6 2. d4 Nf6 3.Nc3 g6 4. f4 Bg7
          5. Bd3 Nc6 6. Nf3 Bg4} 

\chessboard[inverse,showmover=false]
\lfigure

\xskakcomment{There are 3 viable options for white. 7. Be3, 7. e5 and 7.d5}

\subsection{7. Be3}

\newchessgame[id=EarlyDeviation,moveid=6b,
              setfen=\xskakget{nextfen},
              id=EarlyDeviation1,moveid=7w]
\mainline{7. Be3}

\chessboard[inverse,showmover=false]
\lfigure

\xskakcomment{Black now attacks white central pawns with e5}

\mainline{7... e5. 8. dxe5 dxe5 9. }

\subsection{7. e5}

\newchessgame[id=EarlyDeviation,moveid=6b,
              setfen=\xskakget{nextfen},
              id=EarlyDeviation2,moveid=7w]
\mainline{7. e5}

\chessboard[inverse,showmover=false]
\lfigure

\subsection{7. d5}

\newchessgame[id=EarlyDeviation,moveid=6b,
              setfen=\xskakget{nextfen},
              id=EarlyDeviation3,moveid=7w]
\mainline{7. d5}

\chessboard[inverse,showmover=false]
\lfigure
